\documentclass[11pt,a4paper]{article}
\usepackage[T1]{fontenc}
\usepackage[utf8]{inputenc}
\usepackage{lmodern}
\usepackage[margin=25mm,headheight=15pt]{geometry}
\usepackage{amsmath,amssymb,amsthm,mathtools,mathrsfs,aliascnt}
\usepackage{microtype,booktabs,tabularx,array,enumitem}
\usepackage{xcolor,fancyhdr,titlesec,needspace}
\definecolor{Ink}{HTML}{243747}
\definecolor{Accent}{HTML}{566348}
\usepackage[colorlinks=true,linkcolor=Ink,citecolor=Accent,urlcolor=Accent]{hyperref}
\usepackage{xurl}
\usepackage[nameinlink,noabbrev]{cleveref}
\titleformat{\section}{\Large\bfseries\color{Ink}}{\thesection}{.65em}{}
\titleformat{\subsection}{\large\bfseries\color{Accent}}{\thesubsection}{.65em}{}
\pagestyle{fancy}\fancyhf{}
\fancyhead[L]{\small Hahn--Hilbert spectral theory}
\fancyhead[R]{\small Infinitesimal shifts and range defects}
\fancyfoot[C]{\thepage}
\setlength{\parskip}{.26em}
\setlength{\parindent}{1.1em}
\setlength{\emergencystretch}{3em}
\widowpenalty=10000\clubpenalty=10000
\setlist{itemsep=.15em,topsep=.3em,leftmargin=2em}
\setcounter{tocdepth}{2}
\numberwithin{equation}{section}
\newtheorem{theorem}{Theorem}[section]
\newaliascnt{proposition}{theorem}
\newtheorem{proposition}[proposition]{Proposition}
\aliascntresetthe{proposition}
\newaliascnt{lemma}{theorem}
\newtheorem{lemma}[lemma]{Lemma}
\aliascntresetthe{lemma}
\newaliascnt{corollary}{theorem}
\newtheorem{corollary}[corollary]{Corollary}
\aliascntresetthe{corollary}
\theoremstyle{definition}
\newaliascnt{definition}{theorem}
\newtheorem{definition}[definition]{Definition}
\aliascntresetthe{definition}
\newaliascnt{example}{theorem}
\newtheorem{example}[example]{Example}
\aliascntresetthe{example}
\theoremstyle{remark}
\newaliascnt{remark}{theorem}
\newtheorem{remark}[remark]{Remark}
\aliascntresetthe{remark}
\newaliascnt{warning}{theorem}
\newtheorem{warning}[warning]{Warning}
\aliascntresetthe{warning}
\crefname{theorem}{theorem}{theorems}
\Crefname{theorem}{Theorem}{Theorems}
\crefname{proposition}{proposition}{propositions}
\Crefname{proposition}{Proposition}{Propositions}
\crefname{lemma}{lemma}{lemmas}
\Crefname{lemma}{Lemma}{Lemmas}
\crefname{corollary}{corollary}{corollaries}
\Crefname{corollary}{Corollary}{Corollaries}
\crefname{definition}{definition}{definitions}
\Crefname{definition}{Definition}{Definitions}
\crefname{example}{example}{examples}
\Crefname{example}{Example}{Examples}
\crefname{remark}{remark}{remarks}
\Crefname{remark}{Remark}{Remarks}
\crefname{warning}{warning}{warnings}
\Crefname{warning}{Warning}{Warnings}
\newcommand{\C}{\mathbb C}
\newcommand{\R}{\mathbb R}
\newcommand{\Q}{\mathbb Q}
\newcommand{\N}{\mathbb N}
\newcommand{\Z}{\mathbb Z}
\newcommand{\No}{\mathbf{No}}
\newcommand{\K}{{K_\Gamma}}
\newcommand{\KR}{{R_\Gamma}}
\newcommand{\HH}{\mathcal H_\Gamma}
\newcommand{\m}{\mathfrak m_\Gamma}
\newcommand{\Ocal}{\mathcal O_\Gamma}
\newcommand{\B}{\mathcal B(H)}
\newcommand{\Acal}{\mathcal A_\Gamma(H)}
\newcommand{\supp}{\operatorname{supp}}
\newcommand{\st}{\operatorname{st}}
\newcommand{\Ran}{\operatorname{Ran}}
\newcommand{\End}{\operatorname{End}}
\newcommand{\coker}{\operatorname{coker}}
\newcommand{\spec}{\sigma}
\newcommand{\abs}[1]{\lvert#1\rvert}
\newcommand{\norm}[1]{\lVert#1\rVert}
\newcommand{\ip}[2]{\langle#1,#2\rangle}
\newcommand{\commit}{608dd23c0539fce73723843949ee627641c27b30}
\newcolumntype{Y}{>{\raggedright\arraybackslash}X}
\hypersetup{pdftitle={Infinitesimal Spectral Thickening and Closed-Range Defects in Hahn--Hilbert Spaces},pdfauthor={Research manuscript prepared for Vladimir Reshetnikov},pdfsubject={Surcomplex Hahn workspaces, automatic adjoint structure, normal spectra and coercive operators}}
\begin{document}
\hypersetup{pageanchor=false}
\begin{titlepage}
\centering
\vspace*{9mm}
{\large\scshape A research manuscript for the Surreal project\par}
\vspace{13mm}
{\Huge\bfseries\color{Ink} Infinitesimal Spectral\\[.15em] Thickening\par}
\vspace{5mm}
{\LARGE and Closed-Range Defects\\[.15em] in Hahn--Hilbert Spaces\par}
\vspace{8mm}
{\large Automatic adjoint structure, exact normal spectra,\\
and the obstruction to coercive inversion\par}
\vspace{9mm}
{\large September 22, 2026\par}
\vspace{8mm}
\begin{minipage}{.95\textwidth}
\begin{abstract}
\noindent
Let $H$ be an ordinary complex Hilbert space and let $\Gamma$ be a nonzero
set-sized divisible ordered abelian group. We study the vector-valued Hahn
space $H((t^\Gamma))$ over $\C((t^\Gamma))$, with its coefficientwise-convolution
inner product. An automatic-structure theorem identifies every everywhere-defined
adjointable operator, without an assumed continuity or support condition, with
a Hahn series of bounded operators on $H$. For an ordinary bounded normal
operator $T$, its scalar-extended algebraic spectrum is exactly
\[
 \spec_\Gamma(T)=\spec_\C(T)\ \cup\
       \bigl(\spec_\C(T)' + \mathfrak m_\Gamma\bigr),
\]
where the prime denotes accumulation points and $\mathfrak m_\Gamma$ the
infinitesimals. Isolated spectral points do not thicken; accumulation points
acquire whole complex infinitesimal monads, while no new eigenvalues appear.
A general noncommutative perturbation theorem shows that an injective leading
operator retains its entire algebraic range defect under every positive-order
Hahn perturbation.

For $D e_n=n^{-1}e_n$ on $\ell^2$, the positive operator
$D^2+t^{2\eta}I$ is coercive but not onto, and has a cokernel of dimension at
least the continuum over the Hahn field. In a complete rank-one valuation
topology its range is closed and proper with zero orthogonal complement.
We also characterize precisely which constant data admit infinitesimal
resolvents and when the extension of an ordinary bounded operator has a
least field-valued norm bound. All results are proved in a specified
set-sized model, which embeds its scalars into surcomplex numbers through
Conway normal forms.
\end{abstract}
\end{minipage}
\vfill
\begin{minipage}{.95\textwidth}\small
\textbf{Status.} Candidate original theorem package, with explicit proofs and
a documented literature comparison. No named published open problem is claimed
solved. Priority is not certified; the general proofs have not been
independently refereed or checked by a proof assistant. The accompanying
program checks finite algebraic identities, not the infinite-dimensional
theorems. The classical Hilbert, Hahn-support, and surreal normal-form inputs
are identified rather than presented as new.
\end{minipage}
\end{titlepage}
\hypersetup{pageanchor=true}
\pagenumbering{roman}
\tableofcontents
\clearpage
\pagenumbering{arabic}

\section{The problem, the model, and the contribution}\label{sec:intro}

\subsection{Why infinite dimension is the next substantial question}
The repository \texttt{VladimirReshetnikov/Surreal} already contains a substantial
finite-dimensional spectral report: Hermitian diagonalization, singular values,
conditioning, determinantal scales, and ramification of positive matrix roots.
Its scope explicitly restricts matrix dimensions to ordinary finite integers
\cite{RepoSpectral}. Its documentation also distinguishes strong Hahn summation
from convergence in the fine topology of the full surreal class
\cite{RepoMap}. Those are important constraints on what a genuinely additional
chapter should do.

This manuscript asks what happens when an ordinary infinite-dimensional
Hilbert space is retained as the coefficient space, rather than replacing a
finite matrix by a larger finite matrix. Over $\C$, positivity, completeness,
and adjoints support a familiar chain of implications: self-adjoint spectra
are real, coercive operators are invertible, and closed subspaces admit
orthogonal projections. The construction below retains a positive-definite
inner product, all ordinary Hilbert vectors as constants, and a large
adjointable operator algebra. In rank one it is complete. Nevertheless,
several of those implications fail, for an exactly identifiable reason:
\emph{a leading coefficient must itself be an ordinary Hilbert vector}.

The central results are not a list of isolated anomalies. An automatic
operator-structure theorem makes the admissible algebra intrinsic. An
injective-leading-term theorem explains the persistence of range defects.
Together with ordinary normal spectral theory, these give an exact spectral
classification and quantitative counterexamples to coercive inversion.

\subsection{Main statements at a glance}
Throughout, $H\ne0$ is a complex Hilbert space, $\Gamma$ is nonzero, divisible,
ordered, abelian and set-sized, and
\[
 \K=\C((t^\Gamma)),\qquad \HH=H((t^\Gamma)).
\]
The valuation is the least support exponent. A prime on an ordinary compact
set denotes its accumulation set, not a derivative and not an essential
spectrum.

\begin{theorem}[Main theorem package]\label{thm:package}
The following assertions hold in this model.
\begin{enumerate}[label=\textup{(\roman*)}]
\item Every everywhere-defined adjointable $\K$-linear operator on $\HH$
has a unique expansion in $\B((t^\Gamma))$. Its coefficient operators are
ordinary bounded operators and have one common well-ordered support.
\item For an ordinary bounded normal $T$, its coefficientwise extension has
\begin{equation}\label{eq:main-spectrum}
 \spec_\Gamma(T)=\spec_\C(T)\cup
       \bigl(\spec_\C(T)'+\m\bigr),
 \qquad \spec_{p,\Gamma}(T)=\spec_{p,\C}(T).
\end{equation}
Here $\spec_\Gamma$ is the failure-of-bijectivity spectrum; it also equals
the spectrum in the adjointable algebra.
\item If $S:V\to V$ is injective and $E$ is an endomorphism-valued Hahn
series of strictly positive order, then
\[
 \coker(S+E)\cong W((t^\Gamma)),\qquad
 V=\Ran S\oplus W,
\]
noncanonically. No commutation hypothesis is required.
\item For $D e_n=n^{-1}e_n$ and $s=t^\eta$, $\eta>0$, the operator
$C=D^2+s^2I$ on $\ell^2((t^\Gamma))$ is bounded, positive, self-adjoint and
coercive, but not surjective. Its cokernel has $\K$-dimension at least
$2^{\aleph_0}$. In rank one its range is closed and proper, with zero
orthogonal complement.
\item The extension of an ordinary bounded $T$ has a least
$\KR$-valued norm bound if and only if $T$ attains its ordinary Hilbert
operator norm. When the least bound exists, it is the ordinary norm itself.
\end{enumerate}
\end{theorem}

Each assertion is proved below. The result about the least norm bound is
not a claim that these operators are unbounded: all operators in (i) have
field-valued bounds. The failure is the absence of a \emph{least} such bound.
Likewise, a spectral point in (ii) need not be an eigenvalue.

\subsection{What is and is not claimed as new}
The ordinary spectral theorem, Baire category, the closed graph theorem,
Hilbert norm estimates and Hahn support lemmas are classical inputs
\cite{Douglas,Neumann}. The non-Archimedean failure of orthogonal decomposition
is also known in other models, including work of Aguayo and Nova
\cite{AN} and Aguayo, Nova and Shamseddine \cite{ANS}. We do not claim to
have discovered that failure in general.

The proposed contribution is the particular combination of automatic
adjoint structure, the exact accumulation-monad spectrum, the arbitrary
positive-support range-defect formula, and the explicit coercive example
with a continuum-dimensional defect in $H((t^\Gamma))$. In the repository
sections and primary literature inspected, we did not find these statements
in this form. This is a qualified report of the search, not a proof of
priority. A theorem may have predecessors under terminology not located by
that search.

\begin{table}[ht]
\centering\small
\begin{tabularx}{\textwidth}{@{}p{.25\textwidth}YY@{}}
\toprule
Strand & Established input or predecessor & Contribution here\\
\midrule
Finite surcomplex matrices & Repository spectral chapter & Infinite coefficient Hilbert spaces\\
Hahn formal calculus & Neumann support lemma & Automatic global support for adjoints\\
Ordinary normal spectra & Hilbert spectral theorem & Exact infinitesimal thickening formula\\
Formal perturbation & Neumann inversion near the identity & Persistence of arbitrary injective range defects\\
Non-Archimedean orthogonality & Known nonorthomodular examples & A closed coercive range with continuum defect\\
Operator norm bounds & Ordinary norm and attainment & Exact least-bound criterion over the Hahn field\\
\bottomrule
\end{tabularx}
\caption{The novelty boundary: classical tools are not renamed as new theorems.}
\end{table}

\subsection{Relation to the existing non-Archimedean literature}
Aguayo--Nova--Shamseddine work with free Banach spaces of countable type over
the complex Levi--Civita field, represented by weighted $c_0$-type spaces
\cite{ANS}. Ishizuka proves spectral results for self-adjoint compactoid
operators over valued fields with formally real residue field
\cite{Ishizuka}. These are nearby but different settings. The constant copy
of $H$ in our space retains every ordinary square-summable vector, while its
valuation topology is discrete. Our extension of an ordinary compact
operator is not thereby compactoid in the new topology.

The June 2026 preprint of Miabey treats diagonal operators and finite-rank
perturbations in a non-Archimedean Hilbert-space setting \cite{Miabey}.
We use it only to delimit related contemporary work; none of its theorems
is an input to the proofs here. Formal resolvent expansions also have a
classical relationship with pole and Drazin-spectrum questions; see
Boasso \cite{Boasso}. Our normal-operator formula is proved directly, so
no Drazin-theoretic equivalence is needed as an unproved bridge.

\section{Hahn scalar and vector spaces}\label{sec:construction}

\subsection{Support conventions and the finite-convolution principle}
A Hahn series with coefficients in a complex vector space $V$ is a family
\[
 x=\sum_{\gamma\in\Gamma}x_\gamma t^\gamma
\]
whose support $\{\gamma:x_\gamma\ne0\}$ is well ordered in the given order
of $\Gamma$. Write $V((t^\Gamma))$ for the resulting set. For $x\ne0$ put
$v(x)=\min\supp x$, and put $v(0)=+\infty$.

A family $(x_j)_{j\in J}$ is \emph{Hahn summable} when its union of supports
is well ordered and only finitely many $j$ contribute at each exponent.
All unqualified infinite formal sums in this manuscript mean this notion.
It is not convergence of ordinary partial sums in the full surreal fine
topology.

\begin{lemma}[Classical support lemma]\label{lem:support}
Let $A,B\subseteq\Gamma$ be well ordered. Then $A+B$ is well ordered and
each $\gamma$ has only finitely many representations $\gamma=a+b$ with
$a\in A$, $b\in B$. If $P\subseteq\Gamma_{>0}$ is well ordered, the monoid
$P^*$ of all finite sums from $P$, including zero, is well ordered, and for
each $\gamma$ only finitely many ordered words in $P$ have sum $\gamma$.
\end{lemma}

This is the standard Hahn--Neumann support lemma \cite{Neumann}.
We use its full finite-word form, not only the assertion that each fixed
power is meaningful. The latter alone would not justify a geometric series.
For the first finiteness assertion, infinitely many distinct $a$ in a
fixed-sum family yield an increasing subsequence in $A$ and hence a
strictly decreasing subsequence $\gamma-a$ in $B$, a contradiction.
The positive-word assertion is the usual Neumann lemma; it is a stated
classical input here, rather than a new combinatorial claim.

The first assertion licenses scalar multiplication on $V((t^\Gamma))$.
It also licenses products of operator-valued Hahn series, even when the
coefficient algebra is noncommutative. A finite number of supports can be
combined by iteration of the lemma.

\begin{lemma}[Support-certified Neumann inverse]\label{lem:neumann}
Let $\mathscr R$ be a unital complex algebra and let
$E\in\mathscr R((t^\Gamma))$ have $v(E)>0$. Then
\begin{equation}\label{eq:neumann}
 (I+E)^{-1}=\sum_{n\ge0}(-E)^n
\end{equation}
is a well-defined two-sided inverse in $\mathscr R((t^\Gamma))$.
It has support in $(\supp E)^*$.
\end{lemma}
\begin{proof}
Expand each coefficient of $E^n$ as a sum over ordered words in $\supp E$.
\Cref{lem:support} makes the union well ordered and all coefficients finite,
including the sum over $n$. The finite identity
$(I+E)\sum_{n=0}^{N}(-E)^n=I+(-1)^N E^{N+1}$ then telescopes coefficientwise
across all word lengths. Right multiplication telescopes as well. This
argument uses associativity but not commutativity of $\mathscr R$.
\end{proof}

\subsection{The field and its surcomplex realization}
Define
\[
 \KR=\R((t^\Gamma)),\qquad \K=\C((t^\Gamma))=\KR[i].
\]
Order $\KR$ by its first nonzero coefficient and conjugate $\K$
coefficientwise. Let $\Ocal=\{z:v(z)\ge0\}$ and
$\m=\{z:v(z)>0\}$, including zero in both sets. Coefficient extraction at
zero gives $\st:\Ocal\to\C$. Every finite scalar has a unique expression
$z=a+\varepsilon$ with $a\in\C$ and $\varepsilon\in\m$.

Divisibility of $\Gamma$ guarantees positive square roots in $\KR$.
For completeness, if $r=c t^\delta(1+u)>0$, where $c>0$ and $v(u)>0$,
then
\[
 \sqrt r=\sqrt c\,t^{\delta/2}
       \sum_{n\ge0}\binom{1/2}{n}u^n.
\]
The support lemma and the formal binomial identity prove its square is $r$;
the positive root is unique in any ordered field. This elementary square-root
construction is all the real-closed-field structure required here.

If $\Gamma$ is a set-sized ordered subgroup of the surreal additive group,
Conway normal forms give a scalar embedding
\begin{equation}\label{eq:surreal-embedding}
 \sum_\gamma c_\gamma t^\gamma
 \longmapsto \sum_\gamma c_\gamma\omega^{-\gamma}\in\No[i].
\end{equation}
The reversed exponent support is reverse well ordered, exactly the normal-form
condition \cite{Gonshor}. This explains the surcomplex interpretation without
making the vector space a proper class. In particular, $\Gamma=\Q$ and
$t=\omega^{-1}$ suffice for all explicit counterexamples.

\begin{warning}[Three distinct completions]
The construction $H((t^\Gamma))$ is neither the algebraic tensor product
$H\otimes_\C\K$ nor a silently chosen sequence space $\ell^2(\K)$.
Nor does its intrinsic valuation topology equal the topology induced by
embedding its scalar field into the full surreal fine topology. All three
distinctions matter to the results.
\end{warning}

\subsection{The Hahn--Hilbert inner product}
We take ordinary and extended inner products to be linear in the first
variable and conjugate-linear in the second. On $\HH$ define
\begin{equation}\label{eq:inner}
 \ip{x}{y}
 =\sum_{\alpha,\beta}
       \ip{x_\alpha}{y_\beta}_H t^{\alpha+\beta}.
\end{equation}
Each ordinary Hilbert inner product in the coefficient is evaluated
\emph{before} the Hahn convolution. For example, a constant vector
$(1/n)_{n\ge1}\in\ell^2$ has its ordinary finite real squared norm, even
though the separate family of its constant-coordinate squared magnitudes
is not Hahn summable. We are not using that inadmissible family as a Hahn
sum.

\begin{proposition}[Positive-definite scalar extension]\label{prop:inner}
Formula \eqref{eq:inner} defines a positive-definite Hermitian form. Its norm
$\norm{x}=\sqrt{\ip{x}{x}}$ satisfies the Cauchy--Schwarz and triangle
inequalities, and
\begin{equation}\label{eq:normvaluation}
 v(\ip{x}{x})=2v(x),\qquad v(\norm{x})=v(x)\quad(x\ne0).
\end{equation}
If $x=t^\delta(u+\text{positive-order terms})$ with $u\ne0$, then
$\norm{x}=t^\delta(\norm{u}_H+\text{positive-order terms})$.
\end{proposition}
\begin{proof}
The finite-convolution lemma gives well-definedness and sesquilinearity.
The coefficient at $2\delta$, where $\delta=v(x)$, is exactly
$\norm{x_\delta}_H^2>0$; all other contributions have larger exponent.
This proves positivity and the valuation statements. Completing a square,
using $\ip{x-ay}{x-ay}\ge0$ with
$a=\ip{x}{y}/\ip{y}{y}$, gives
$\abs{\ip{x}{y}}^2\le\ip{x}{x}\ip{y}{y}$. The usual expansion of
$\norm{x+y}^2$, followed by this inequality, gives the triangle inequality.
These are finite algebraic arguments over the ordered field $\KR$.
\end{proof}

\section{Every adjoint has a global Hahn expansion}\label{sec:adjoints}

\subsection{The apparently smaller coefficient algebra}
Let
\[
 \Acal=\B((t^\Gamma)).
\]
Thus an element $A=\sum_\gamma A_\gamma t^\gamma$ has bounded ordinary
operator coefficients and one well-ordered support. There is no uniform
ordinary bound on the norms of all its coefficients.
It acts by
\begin{equation}\label{eq:operator-convolution}
 (Ax)_\delta=\sum_{\gamma+\alpha=\delta}A_\gamma x_\alpha.
\end{equation}
The support lemma proves this action is defined on every $x\in\HH$.
The series $A^*=\sum_\gamma A_\gamma^*t^\gamma$ satisfies
$\ip{Ax}{y}=\ip{x}{A^*y}$, by finite coefficientwise rearrangement.

One might suspect that $\Acal$ is only a convenient restricted operator
class. The next theorem says that the existence of an everywhere-defined
adjoint forces exactly this class, without any advance assumption that an
operator preserves Hahn sums.

\begin{theorem}[Automatic adjoint structure]\label{thm:adjoint}
Let $A,B:\HH\to\HH$ be everywhere-defined $\K$-linear maps such that
\begin{equation}\label{eq:adjoint-condition}
 \ip{Ax}{y}=\ip{x}{By}\qquad(x,y\in\HH).
\end{equation}
There are unique bounded operators $A_\gamma\in\B$ with well-ordered
nonzero support such that $A$ acts by \eqref{eq:operator-convolution}.
Moreover, $B_\gamma=A_\gamma^*$. Consequently,
\[
 \End_{\mathrm{adj}}(\HH)\cong\B((t^\Gamma))
\]
as $*$-algebras. No separability of $H$ is assumed.
\end{theorem}
\begin{proof}
\emph{Step 1: bounded coefficient operators.}
For constant vectors $u,w\in H$, define
$A_\gamma u=[t^\gamma]Au$ and $B_\gamma w=[t^\gamma]Bw$.
Extracting the coefficient at $\gamma$ in \eqref{eq:adjoint-condition} gives
\begin{equation}\label{eq:ordinaryadj}
 \ip{A_\gamma u}{w}_H=\ip{u}{B_\gamma w}_H.
\end{equation}
Both maps are defined on all of $H$. To see that $A_\gamma$ has closed
graph, let $u_n\to u$ and $A_\gamma u_n\to z$ in the ordinary Hilbert
norm. Testing against any fixed $w$ in \eqref{eq:ordinaryadj} gives
$\ip{z}{w}=\ip{u}{B_\gamma w}=\ip{A_\gamma u}{w}$.
Hence $z=A_\gamma u$. The ordinary closed graph theorem makes $A_\gamma$
bounded. The same argument applies to $B_\gamma$, and
\eqref{eq:ordinaryadj} identifies the two as adjoints.

\emph{Step 2: one well-ordered support.}
Let $S=\{\gamma:A_\gamma\ne0\}$. If $S$ were not well ordered, ordinary
choice permits a strictly decreasing sequence
$\gamma_1>\gamma_2>\cdots$ in $S$. Each $\ker A_{\gamma_n}$ is a proper
closed subspace of $H$, and hence nowhere dense. Baire's theorem gives
$u\in H$ outside their countable union. Then the support of $Au$ contains
every $\gamma_n$, contradicting its being a Hahn series. Thus $S$ is well
ordered. Because an ordinary bounded operator is zero exactly when its
adjoint is zero, the coefficient family of $B$ has the same support.

\emph{Step 3: constants determine the full operator.}
Let $\widetilde A,\widetilde B$ be the convolution operators built from
these coefficient families. They are adjoints by the support lemma and agree
with $A,B$ on constants. For arbitrary $x\in\HH$ and constant $w\in H$,
\[
 \ip{Ax}{w}=\ip{x}{Bw}
 =\ip{x}{\widetilde B w}=\ip{\widetilde A x}{w}.
\]
If a Hahn vector pairs to zero with every constant $w$, then each of its
ordinary coefficients pairs to zero with all of $H$, so every coefficient
vanishes. Hence $Ax=\widetilde A x$. The analogous conclusion holds for $B$.
Uniqueness follows by evaluating on constant vectors and extracting
coefficients. Composition and the involution are the finite Hahn convolution
operations, proving the algebra assertion.
\end{proof}

\begin{corollary}[Automatic summability and boundedness]\label{cor:autobounded}
Every everywhere-defined adjointable operator preserves Hahn-summable
families. It also admits a bound
\begin{equation}\label{eq:orderbound}
 \norm{Ax}\le M\norm{x}\qquad(x\in\HH)
\end{equation}
for some $M\in\KR_{\ge0}$. If $A\ne0$ and
$\gamma_0=\min\supp A$, any ordinary real $c>\norm{A_{\gamma_0}}_H$
gives the valid bound $M=c t^{\gamma_0}$.
\end{corollary}
\begin{proof}
For a Hahn-summable family, finite-convolution finiteness with $\supp A$
allows the two sums to be interchanged, and their union of supports remains
well ordered. For the bound, write the leading term of $x$ as
$t^\delta u$. At exponent $2(\gamma_0+\delta)$, the coefficient of
$M^2\norm{x}^2-\norm{Ax}^2$ is
$c^2\norm{u}_H^2-\norm{A_{\gamma_0}u}_H^2>0$.
This remains true when $A_{\gamma_0}u=0$, since then the leading output
moves to a higher exponent. Thus the difference is positive for nonzero $x$.
\end{proof}

\begin{remark}
The Baire argument is a uniform-support theorem, not a uniform bound on
$\norm{A_\gamma}_H$. The proof also explains why unrestricted ordinary
algebraic endomorphisms used later for choosing range complements are not
automatically adjointable.
\end{remark}

\section{An injective leading term preserves the whole range defect}\label{sec:defect}

\subsection{A theorem independent of Hilbert structure}
The next result applies to any complex vector space $V$. Write
$V_\Gamma=V((t^\Gamma))$. An ordinary linear map on $V$ extends
coefficientwise, even if $V$ also carries a topology and the map is
discontinuous in that topology.

\begin{theorem}[Positive-order defect rigidity]\label{thm:defect}
Let $S:V\to V$ be injective. Let
$E\in\End_\C(V)((t^\Gamma))$ have strictly positive order, or be zero,
and put $A=S+E$. Choose an algebraic complement
\begin{equation}\label{eq:complement}
 V=\Ran S\oplus W.
\end{equation}
Then
\begin{equation}\label{eq:defectdecomp}
 V_\Gamma=A(V_\Gamma)\oplus W((t^\Gamma)).
\end{equation}
In particular, $A$ is injective, $v(Ax)=v(x)$ for $x\ne0$, and
$\coker A\cong W((t^\Gamma))$ as $\K$-vector spaces. The isomorphism depends
on the chosen complement, but the nonvanishing of the defect does not.
\end{theorem}
\begin{proof}
At the least exponent $\delta$ of a nonzero $x$, the coefficient of $Ax$
is $Sx_\delta\ne0$. All contributions from $E$ have greater exponent.
This proves injectivity and the valuation assertion.

Define $U:V\to V$ by
$U(Sv+w)=v$ for $v\in V$, $w\in W$. Then $US=I$ and $\ker U=W$.
Extend $U$ coefficientwise and use \cref{lem:neumann} to set
\[
 R=(I+UE)^{-1},\qquad L=RU,\qquad \Pi=I-AL.
\]
The coefficients of $U$ need not be bounded in any additional topology;
all these formulas are legitimate in the ordinary algebraic endomorphism
Hahn algebra. The identities
\[
 LA=R(US+UE)=R(I+UE)=I
\]
and
\[
 U\Pi=U-UA R U=U-(I+UE)R U=0
\]
show that $\Pi$ kills the range of $A$ and takes values in
$\ker U=W((t^\Gamma))$. On that latter space $L$ vanishes, so $\Pi$ is the
identity. Finally, every $x$ decomposes as $ALx+\Pi x$; the intersection
of the two summands is zero because $LA=I$ and $L|_{W((t^\Gamma))}=0$.
\end{proof}

\begin{corollary}[No positive-order repair]\label{cor:norepair}
An injective nonsurjective ordinary operator cannot become surjective on
$V_\Gamma$ by adding any globally supported positive-order Hahn
perturbation. This includes perturbations that do not commute with it and
perturbations with infinitely many positive exponents.
\end{corollary}

\begin{corollary}[Explicit independent obstruction classes]\label{cor:independent}
Suppose $(f_j)_{j\in J}\subset V$ have linearly independent classes in
$V/\Ran S$ over $\C$. Their constant Hahn extensions have linearly
independent classes in $V_\Gamma/(S+E)V_\Gamma$ over $\K$. Consequently,
\[
 \dim_\K\coker(S+E)\ge\dim_\C(V/\Ran S).
\]
\end{corollary}
\begin{proof}
In a nonzero finite sum $f=\sum_j c_j f_j$ with $c_j\in\K$, take the least
exponent among the nonzero $c_j$. The coefficient at that exponent is a
nonzero finite complex combination of the $f_j$, and is not in $\Ran S$.
But the leading coefficient of $(S+E)x$ is always in $\Ran S$.
\end{proof}

\begin{warning}
The theorem gives the actual vector-space quotient, not only a first-order
approximation to it. It does \emph{not} assert
$\dim_\K W((t^\Gamma))=\dim_\C W$ for arbitrary infinite-dimensional $W$.
The safe dimension statement is the lower bound just proved.
\end{warning}

\subsection{Why the result is more informative than a determinant test}
There is no finite determinant for an arbitrary infinite-dimensional $S$.
An entrywise infinitesimal perturbation may even make every finite section
invertible, while the leading ordinary coefficient of a proposed inverse
vector is excluded from $H$. The exact splitting \eqref{eq:defectdecomp}
shows that this is not an accidental choice of data: the full algebraic
cokernel persists. On the other hand, if $S$ is an ordinary bijection,
$S+E=S(I+S^{-1}E)$ is invertible by the same support-certified Neumann
formula. Thus the injective case has a sharp dichotomy, not merely a
negative example.

\section{The exact spectrum of an ordinary normal operator}\label{sec:spectrum}

\subsection{Two spectral definitions and one ordinary reduction}
For a $\K$-linear endomorphism $A$ of $\HH$, define
\[
 \spec_\Gamma^{\mathrm{alg}}(A)
 =\{\lambda\in\K:A-\lambda I\text{ is not bijective}\}.
\]
For $A\in\Acal$, define $\spec_\Gamma^{\mathrm{adj}}(A)$ by failure of
invertibility in the unital algebra $\Acal$. These definitions need not be
identified in advance. In the theorem below, explicit inverses and explicit
range obstructions will prove their equality.

An ordinary bounded operator $T$ is regarded as a constant member of
$\Acal$. We often suppress the extension subscript on $T$.

\begin{lemma}[Normal isolated-point decomposition]\label{lem:normalgap}
Let $T\in\B$ be normal and $a\in\spec_\C(T)$. Put
$N=\ker(T-aI)$, $M=N^\perp$, and $S=(T-aI)|_M$. Then $M,N$ reduce $T$
and $S$ is injective. If $a$ is isolated in $\spec_\C(T)$, then $N\ne0$
and $S$ is boundedly invertible on $M$. If $a$ is an accumulation point,
then $S$ is not surjective.
\end{lemma}
\begin{proof}
These facts follow from the ordinary normal spectral theorem
\cite{Douglas}; here is the precise use of it. The spectral projection of
$\{a\}$ is the orthogonal projection onto $N$. For isolated $a$, the
function $(z-a)^{-1}$ on the rest of the compact spectrum is bounded, so
spectral functional calculus gives an inverse on $M$. The projection cannot
be zero, since otherwise $a$ would not belong to the spectrum.
For accumulation $a$, if $S$ were onto, its injectivity and the bounded
inverse theorem would give a bounded inverse. The spectrum of $T$ would
then be contained in $\{a\}$ together with a compact set separated from
$a$, contradicting accumulation. This also covers $N=0$.
\end{proof}

\subsection{Classification and all inverse formulas}
\begin{theorem}[Accumulation-monad spectrum]\label{thm:spectrum}
For every ordinary bounded normal $T\in\B$,
\begin{equation}\label{eq:exactspectrum}
 \spec_\Gamma^{\mathrm{alg}}(T)
 =\spec_\Gamma^{\mathrm{adj}}(T)
 =\spec_\C(T)\cup\bigl(\spec_\C(T)'+\m\bigr).
\end{equation}
No scalar of negative valuation is in this spectrum. A nonzero
infinitesimal displacement of an isolated ordinary spectral point is in
the resolvent, whereas every infinitesimal displacement of an accumulation
point remains in the spectrum.
\end{theorem}
\begin{proof}
We partition all scalars $\lambda\in\K$ into exhaustive cases.

\emph{Infinite scalar.} If $v(\lambda)<0$, then
\begin{equation}\label{eq:infiniteres}
 (T-\lambda I)^{-1}
 =-\sum_{n\ge0}\lambda^{-n-1}T^n.
\end{equation}
This is a support-certified Neumann inverse: the scalar $\lambda^{-1}$
has positive valuation. It belongs to $\Acal$, irrespective of the
ordinary growth of $\norm{T^n}$.

\emph{Finite scalar with nonspectral residue.} Write
$\lambda=a+\varepsilon$, where $a\in\C$, $\varepsilon\in\m$.
If $a\notin\spec_\C(T)$, let $B=(T-aI)^{-1}\in\B$. Then
\begin{equation}\label{eq:regularres}
 (T-aI-\varepsilon I)^{-1}
 =\sum_{n\ge0}\varepsilon^n B^{n+1}.
\end{equation}
The formula includes $\varepsilon=0$. For nonzero $\varepsilon$, use
\cref{lem:neumann} on $\varepsilon B$; coefficients need not satisfy
any ordinary norm-convergence estimate.

\emph{Undisplaced ordinary spectral point.} If $\varepsilon=0$ and
$a\in\spec_\C(T)$, then $T-aI$ is not an ordinary bijection, by the bounded
inverse theorem. A nonzero ordinary kernel vector remains a Hahn kernel
vector. If there is no kernel but surjectivity fails, choose a constant
$f$ outside its ordinary range. The coefficient at zero in any putative
solution would have to solve $(T-aI)y_0=f$, which is impossible. Thus the
Hahn operator is not bijective.

\emph{Displaced isolated point.} Suppose $a$ is isolated and
$\varepsilon\ne0$. Use \cref{lem:normalgap}; let $P,Q$ project onto $N,M$.
The inverse is
\begin{equation}\label{eq:isolatedres}
 (T-aI-\varepsilon I)^{-1}
 =-\varepsilon^{-1}P+
   \sum_{n\ge0}\varepsilon^n S^{-n-1}Q.
\end{equation}
The first term is a single scalar Hahn series times a bounded projection;
the second is a positive-order Neumann series on $M$. Both have
well-ordered operator support and lie in $\Acal$. The finite orthogonal
splitting verifies the inverse on both summands. If $M=0$, the second
summand is simply zero.

\emph{Displaced accumulation point.} Suppose $a\in\spec_\C(T)'$ and
$\varepsilon\ne0$. On $M$ in \cref{lem:normalgap}, the ordinary $S$ is
injective and not onto. The perturbation $-\varepsilon I$ has strictly
positive order. By \cref{thm:defect}, $S-\varepsilon I$ is not onto
$M((t^\Gamma))$. The constant decomposition $H=N\oplus M$ extends to
a decomposition of Hahn spaces, and the full operator is block diagonal.
It is therefore not onto either. Explicitly, a constant
$f\in M\setminus\Ran S$ cannot be its $M$-component output.

These cases prove exactly the stated set formula. In every nonspectral
case we constructed an inverse in $\Acal$; in every spectral case we
proved failure of algebraic bijectivity. This proves equality of the two
spectral definitions as well.
\end{proof}

\subsection{Point spectrum is rigid, despite spectral thickening}
\begin{theorem}[No new normal eigenvalues]\label{thm:point}
For ordinary bounded normal $T$,
\[
 \spec_{p,\Gamma}(T)=\spec_{p,\C}(T),\qquad
 \ker(T-aI)=\ker_H(T-aI)((t^\Gamma))\quad(a\in\C).
\]
In particular, the new nonreal spectral points of a self-adjoint operator
are not eigenvalues.
\end{theorem}
\begin{proof}
Infinite scalars are in the resolvent by \cref{thm:spectrum}. Suppose
$(T-aI-\varepsilon I)x=0$ with $x\ne0$, $\varepsilon\in\m$.
At its leading exponent $\delta$, one obtains
$(T-aI)x_\delta=0$. Thus $a$ is an ordinary eigenvalue and
$x_\delta\in N=\ker_H(T-aI)$.
If $\varepsilon\ne0$, apply the ordinary orthogonal projection $P$ onto
$N$. Normality implies $P(T-aI)=0$, so the full equation gives
$-\varepsilon Px=0$. Hence $Px=0$, contradicting
$Px_\delta=x_\delta\ne0$. Therefore $\varepsilon=0$.
For a constant $a$, the kernel equation is coefficientwise, yielding the
claimed eigenspace identity.
\end{proof}

\begin{example}[Normality is indispensable]\label{ex:backshift}
Let $B$ be the backward shift on $\ell^2(\N_0)$:
$Be_0=0$, $Be_n=e_{n-1}$ for $n\ge1$. For $s=t^\eta$,
\[
 x=\sum_{n\ge0}s^n e_n\in\HH,\qquad Bx=sx.
\]
The support $\{n\eta:n\ge0\}$ is well ordered. This produces a new
nonconstant eigenvalue for the nonnormal operator $B$, in direct contrast
to \cref{thm:point}. No ordinary Hilbert convergence of the vector series
is asserted or needed.
\end{example}

\section{Consequences that separate finite and infinite spectral geometry}\label{sec:consequences}

\subsection{Self-adjoint spectra are real exactly in the finite-spectrum case}
\begin{corollary}[A sharp spectral-reality criterion]\label{cor:reality}
Let $T$ be ordinary bounded self-adjoint. The following are equivalent:
\begin{enumerate}[label=\textup{(\roman*)}]
\item $\spec_\Gamma(T)\subseteq\KR$;
\item $\spec_\C(T)$ is finite;
\item $T$ has a finite ordinary orthogonal spectral resolution
$T=\sum_{j=1}^{r}a_jP_j$ with $a_j\in\R$.
\end{enumerate}
There is no finite-rank hypothesis on the projections $P_j$.
\end{corollary}
\begin{proof}
The usual self-adjoint spectrum is a nonempty compact subset of $\R$.
A finite set has no accumulation points, so \cref{thm:spectrum} proves
(ii)$\Rightarrow$(i). An infinite compact subset has an accumulation point
$a$. For any $s=t^\eta$, $\eta>0$, the scalar $a+is$ belongs to
$\spec_\Gamma(T)$ and does not belong to $\KR$, proving the converse.
The equivalence with (iii) is the ordinary spectral theorem applied to a
finite spectral set.
\end{proof}

The criterion is finite \emph{spectrum}, not finite dimension. For example,
an infinite-rank orthogonal projection with infinite-dimensional kernel has
spectrum exactly $\{0,1\}$ even after Hahn extension. Infinite multiplicity
at an isolated point does not create a monad. The accumulation set in
\eqref{eq:exactspectrum} therefore cannot be replaced by an essential
spectrum that includes such points.

\subsection{Ordinary compact operators and a continuous-spectrum example}
\begin{corollary}[Ordinary compact normal extension]\label{cor:compact}
If $T$ is an ordinary compact normal operator of infinite rank, then
\[
 \spec_\Gamma(T)=\m\ \cup\
        \{\text{nonzero ordinary eigenvalues of }T\}.
\]
If its ordinary spectrum is finite, there is no thickening.
\end{corollary}
\begin{proof}
The ordinary compact normal spectral theorem says that nonzero spectral
points are isolated eigenvalues of finite multiplicity and zero is the
only possible accumulation point. Infinite rank forces zero to be an
accumulation point. Apply \cref{thm:spectrum}.
\end{proof}

Here ``compact'' describes the original operator on the ordinary Hilbert
space. It makes no claim of compactness or compactoidness in the Hahn
valuation topology. The sequence $(De_n)$ for $De_n=n^{-1}e_n$, for
instance, has pairwise valuation distance one, despite $\norm{De_n}_H\to0$
in the ordinary norm.

For another concrete example, let $T$ be multiplication by the coordinate
on $L^2([0,1],\C)$. Its ordinary spectrum is $[0,1]$, it has no eigenvalues,
and every point of its spectrum is an accumulation point. Thus
\[
 \spec_\Gamma(T)=[0,1]+\m,\qquad \spec_{p,\Gamma}(T)=\varnothing.
\]
The usual facts here can be seen directly: away from $[0,1]$ the reciprocal
multiplier is bounded; near any point of $[0,1]$, vectors supported in
arbitrarily small intervals prevent a bounded inverse; a singleton has
measure zero and cannot support a nonzero eigenvector. This gives full
infinitesimal fibers above a continuous real interval, without adding a
single eigenvector.

\subsection{Enlarging the scalar workspace does not repair the defect}
\begin{corollary}[Persistence under ordered-group extension]\label{cor:extension}
Let $\Gamma\subseteq\Delta$ be ordered divisible groups and use the natural
coefficient-preserving embedding $K_\Gamma\subseteq K_\Delta$. For constant
normal $T$,
\[
 \spec_\Delta(T)\cap K_\Gamma=\spec_\Gamma(T).
\]
Moreover, the nonsurjectivity in \cref{cor:norepair} persists after such an
extension for the same $S$ and $E$.
\end{corollary}
\begin{proof}
An element of $K_\Gamma$ has the same leading valuation, residue and
infinitesimal status in the larger group. Apply the exact spectral formula.
For the second assertion use the same complement $W$ in
\cref{thm:defect}, now forming $W((t^\Delta))$.
\end{proof}

This is a genuine surcomplex-workspace statement. Adding further surreal
scales does not turn a prohibited leading ordinary Hilbert coefficient into
an allowed one. It does not assert that an unspecified larger class-sized
notion of vector could never behave differently.

\section{Which constant data have an infinitesimal resolvent?}\label{sec:resolventdata}

\subsection{The inverse-smooth domain}
Even when $S-sI$ is not onto, some constant right-hand sides have a unique
solution. The exact criterion is stronger than belonging merely to
$\Ran S$.

\begin{theorem}[All inverse moments are necessary and sufficient]\label{thm:smooth}
Let $S:V\to V$ be injective and $s=t^\eta$ with $\eta>0$. For constant
$f\in V$, the equation
\begin{equation}\label{eq:constres}
 (S-sI)y=f,\qquad y\in V((t^\Gamma)),
\end{equation}
has a solution if and only if
\begin{equation}\label{eq:inversesmooth}
 f\in\bigcap_{k\ge1}\Ran S^k.
\end{equation}
In that case the solution is unique and is
\begin{equation}\label{eq:smoothsolution}
 y=\sum_{n\ge0}s^n S^{-n-1}f.
\end{equation}
The inverse powers denote the unique preimages, not an everywhere-defined
inverse operator on $V$.
\end{theorem}
\begin{proof}
Sufficiency follows by coefficientwise telescoping of
\eqref{eq:smoothsolution}; its support is contained in the well-ordered set
$\{n\eta:n\ge0\}$. Injectivity follows from \cref{thm:defect}.

For necessity, partition the support of a proposed solution according to
cosets of $\Z\eta$ in $\Gamma$. Restricting a Hahn series to a coset preserves
well-ordering, and both $S$ and multiplication by $s$ preserve that coset.
On a coset other than $\Z\eta$, the restricted equation is homogeneous.
Its least nonzero coefficient would be killed by $S$, which is impossible.
Thus the support lies in $\Z\eta$.

A well-ordered subset of $\Z\eta$ is bounded below in its integer index.
For $f\ne0$, the valuation-preservation assertion of \cref{thm:defect}
forces the least index to be zero. Write $y=\sum_{n\ge0}s^n y_n$, allowing
zero coefficients. The equation gives
\[
 Sy_0=f,\qquad Sy_n=y_{n-1}\quad(n\ge1).
\]
Consequently $S^{n+1}y_n=f$ for every $n$, proving
\eqref{eq:inversesmooth} and the formula. The zero-data case is immediate.
\end{proof}

\begin{example}[The diagonal inverse-smooth space]\label{ex:smoothD}
For $De_n=n^{-1}e_n$ on $H=\ell^2(\N_{\ge1})$, the admissible constant data
are exactly
\begin{equation}\label{eq:rapidspace}
 H_\infty=\left\{f=(f_n):
       \sum_{n\ge1}n^{2k}|f_n|^2<\infty
       \text{ for every integer }k\ge1\right\}.
\end{equation}
The formula follows because $D^{-k}f=(n^k f_n)$. Finitely supported vectors
belong to $H_\infty$, so this subspace is dense in the \emph{ordinary}
Hilbert topology. Yet $f_n=n^{-2}$ does not belong to it, although
$f\in\Ran D$. Thus even some data with an ordinary first preimage lack an
infinitesimal resolvent.
\end{example}

The same statement for $D-isI$ follows by replacing the recurrence with
$Dy_n=i y_{n-1}$. Its solution, when it exists, is
$\sum_{n\ge0}i^n s^nD^{-n-1}f$. No uniform bound on the ordinary norms of
these coefficients is required, only that each one actually belongs to $H$.

\subsection{Why individual scalar inverses are insufficient}
Every diagonal scalar $n^{-1}-s$ is invertible in $\K$ and has expansion
\[
 (n^{-1}-s)^{-1}=\sum_{k\ge0}n^{k+1}s^k.
\]
But applying these inverses coordinatewise to $f$ requires the entire
coefficient vector $(n^{k+1}f_n)_n$ to lie in $\ell^2$ at each $k$.
The Hahn vector space imposes this condition before any comparison of
higher scales. There is therefore no contradiction between invertibility
of every diagonal entry and failure of surjectivity of the infinite
operator. A putative coordinatewise inverse has left the chosen vector
space, rather than found a missing scalar inverse.

\section{Coercive operators with continuum-dimensional defects}\label{sec:coercive}

\subsection{An explicit large ordinary cokernel}
Continue with $H=\ell^2(\N_{\ge1})$ and $De_n=n^{-1}e_n$.
The operator $D$ is bounded, positive, self-adjoint and injective. It is not
onto because $(n^{-1})_n\in\ell^2$ has no preimage in $\ell^2$.

\begin{lemma}[Continuum many independent defect vectors]\label{lem:continuum}
For $r\in(1/2,3/2]$, let $f^{(r)}=(n^{-r})_{n\ge1}\in H$.
Their classes in $H/\Ran D$ are linearly independent over $\C$.
\end{lemma}
\begin{proof}
Take distinct $r_1<\cdots<r_k$ and complex coefficients with $c_1\ne0$.
For $g=\sum_{j=1}^k c_j f^{(r_j)}$, one has
\[
 n g_n=c_1 n^{1-r_1}\bigl(1+o(1)\bigr).
\]
Its squared magnitude is bounded below by a positive real multiple of
$n^{2-2r_1}$ for sufficiently large $n$. Since $r_1\le3/2$, the exponent is
at least $-1$, and the series of those squared magnitudes diverges.
Thus $(n g_n)\notin\ell^2$, which is exactly $g\notin\Ran D$.
There are continuum many choices of $r$.
\end{proof}

\begin{theorem}[Every positive-order perturbation retains a large defect]\label{thm:continuum}
For every $E\in\B((t^\Gamma))$ of positive order, including zero,
\[
 \dim_\K\coker(D+E)\ge2^{\aleph_0}.
\]
The same lower bound holds with leading operator $D^2$.
\end{theorem}
\begin{proof}
Apply \cref{cor:independent} and \cref{lem:continuum}. Since
$\Ran D^2\subseteq\Ran D$, the same vectors remain independent modulo
$\Ran D^2$.
\end{proof}

\subsection{A nonreal spectral point with a strong lower bound}
For $s=t^\eta$, $\eta>0$, set $A=D-isI$. Since $D=D^*$,
\begin{equation}\label{eq:lowerA}
 A^*A=D^2+s^2I,\qquad
 \norm{Ax}^2=\norm{Dx}^2+s^2\norm{x}^2
             \ge s^2\norm{x}^2.
\end{equation}
Nevertheless, $A$ is not onto by \cref{thm:defect}. Thus the usual lower-bound
argument for excluding nonreal self-adjoint spectrum stops precisely at
surjectivity, not at injectivity.

The operator even has an attained upper norm bound:
\begin{equation}\label{eq:normA}
 \norm{Ax}\le\sqrt{1+s^2}\,\norm{x},\qquad
 \norm{Ae_1}=\sqrt{1+s^2}.
\end{equation}
Indeed $D$ is an ordinary contraction and its extension remains a contraction;
for this diagonal operator the inequality also follows by writing
$I-D^2=R^*R$ with ordinary bounded diagonal
$Re_n=\sqrt{1-n^{-2}}e_n$. Equations \eqref{eq:lowerA}--\eqref{eq:normA}
show that neither lack of a bound nor lack of an attained norm explains the
failure of inversion.

\subsection{Positive coercivity does not imply surjectivity}
\begin{theorem}[A coercive noninvertible self-adjoint operator]\label{thm:coercive}
Let
\begin{equation}\label{eq:C}
 C=D^2+s^2I.
\end{equation}
Then $C$ is everywhere-defined, adjointable and self-adjoint, with
\begin{align}
 \ip{Cx}{x}&\ge s^2\ip{x}{x},\label{eq:coercive}\\
 \norm{Cx}&\le(1+s^2)\norm{x}.\label{eq:Cbound}
\end{align}
The upper bound is attained at $e_1$. The operator is injective but not
surjective, and its cokernel has dimension at least $2^{\aleph_0}$ over $\K$.
Specifically, the constant vector $f=(n^{-2})_{n\ge1}$ is not in its range.
\end{theorem}
\begin{proof}
Self-adjointness and \eqref{eq:coercive} follow from
$\ip{D^2x}{x}=\ip{Dx}{Dx}\ge0$.
The contraction property of $D^2$, the triangle inequality and
$Ce_1=(1+s^2)e_1$ give \eqref{eq:Cbound} and its attainment.
The leading coefficient operator $D^2$ is injective, so
\cref{thm:defect} gives injectivity, and \cref{thm:continuum} gives the
cokernel bound.

For a direct witness, if $Cy=f$, valuation preservation forces $v(y)=0$.
The coefficient at zero must solve $D^2y_0=f$, that is,
$(y_0)_n=1$ for every $n$. This is not an $\ell^2$ vector. Thus no such
$y$ exists.
\end{proof}

For clarity, coercivity is a bound by a strictly positive element $s^2$ of
the ordered scalar field. It is not a bound by a positive \emph{ordinary
real} independent of the infinitesimal scale. If instead
$C\ge cI$ with ordinary $c>0$ and an ordinary bounded positive leading
operator with this bound, the classical inverse and the positive-order
Neumann construction would apply. The infinitesimal scale is exactly where
the distinction becomes mathematically significant.

\begin{corollary}[Failure of a specified Lax--Milgram analogue]\label{cor:lax}
On this Hahn inner-product space, a bounded coercive sesquilinear form need
not solve all variational problems, even when the right-hand side is
represented by an actual vector. In particular,
\[
 b(u,v)=\ip{Cu}{v},\qquad \ell(v)=\ip{f}{v},
 \qquad f=(n^{-2})_n,
\]
defines a bounded coercive form and a bounded conjugate-linear functional,
but there is no $u$ satisfying $b(u,v)=\ell(v)$ for all $v\in\HH$.
\end{corollary}
\begin{proof}
Cauchy--Schwarz and \eqref{eq:Cbound} bound the form; \eqref{eq:coercive}
is coercivity, and $\abs{\ell(v)}\le\norm{f}\norm{v}$.
If the identity held for all $v$, nondegeneracy of the inner product would
give $Cu=f$, contradicting \cref{thm:coercive}.
\end{proof}

This does not contradict the ordinary complex Lax--Milgram theorem: the
scalar field, the topology and the duality structure have changed. It is
an explicit obstruction to transferring that theorem using only its
familiar-looking algebraic inequalities.

\section{Completeness does not restore orthogonal projection}\label{sec:topology}

\subsection{The rank-one topology and its completeness}
In this section only, assume $\Gamma$ is a nonzero divisible ordered subgroup
of $\R$. For $x\ne0$ define
\[
 \rho(x)=e^{-v(x)},\qquad \rho(0)=0,\qquad d(x,y)=\rho(x-y).
\]
The scalar absolute value here is $|a|_v=e^{-v(a)}$, not the
$\KR$-valued modulus. The metric is ultrametric. Its topology agrees with
the topology of the field-valued norm balls on $\HH$: a sufficiently
higher valuation makes a vector smaller than any prescribed positive
radius, and choosing a still higher monomial radius gives the converse
inclusion of neighborhood bases.

It is not the ordinary Hilbert topology on the constant copy of $H$.
Every nonzero constant vector has $\rho$-size one. Nor is it the fine
subspace topology from the full surcomplex class, where the set-sized
scalar workspace is discrete \cite{RepoMap}.

\begin{proposition}[Rank-one valuation completeness]\label{prop:complete}
For any complex vector space $V$, the metric space $V((t^\Gamma))$ with
metric $d$ is complete. In particular, $\HH$ is complete in this topology.
\end{proposition}
\begin{proof}
Let $(x_n)$ be Cauchy. For each real cutoff $R$, all sufficiently late
$x_n$ have exactly the same coefficients at exponents $\gamma\le R$,
since eventually $v(x_n-x_m)>R$. Define $x_\gamma$ to be the eventually
stable coefficient.

The support of this coefficient family is well ordered. Otherwise it would
contain a strictly decreasing sequence $\gamma_1>\gamma_2>\cdots$.
Choose a cutoff at least $\gamma_1$. Below that cutoff the entire stable
family agrees with one sufficiently late Hahn vector $x_N$, whose
well-ordered support cannot contain that sequence. This is a contradiction.
Thus $x=\sum x_\gamma t^\gamma$ belongs to $V((t^\Gamma))$.
By the same stabilization statement, $v(x_n-x)>R$ eventually for every
real $R$, proving convergence. Notice that no coefficient-space topology
was used: stabilization is exact, not approximate.
\end{proof}

\subsection{A proper closed subspace with zero orthogonal complement}
\begin{theorem}[Closed-range orthogonality failure]\label{thm:closedrange}
For the coercive $C=D^2+s^2I$, its range
$M=C\HH$ is closed and proper in the complete rank-one space $\HH$, but
\begin{equation}\label{eq:orthdefect}
 M^\perp=\{0\},\qquad M^{\perp\perp}=\HH\ne M.
\end{equation}
In particular, $M$ is not the range of a self-adjoint idempotent.
\end{theorem}
\begin{proof}
By \cref{thm:defect}, $v(Cx)=v(x)$; hence $C$ is an isometry for $\rho$.
The image of a complete metric space under an isometry is complete and
therefore closed in any metric ambient space. It is proper by
\cref{thm:coercive}.

If $z\in M^\perp$, then $\ip{Cx}{z}=0$ for every $x$. Since $C=C^*$,
$\ip{x}{Cz}=0$ for every $x$, so $Cz=0$. Injectivity gives $z=0$.
This proves \eqref{eq:orthdefect}. A self-adjoint idempotent $P$ has
$\HH=\Ran P\oplus\ker P$ and
$\ker P=(\Ran P)^\perp$, by the finite algebraic projection identities.
If its range were $M$, the complement would be zero and the range all of
$\HH$, a contradiction.
\end{proof}

The displayed range is not orthogonally closed, although it is metrically
closed. Thus the usual Hilbert identification of these two closure notions
fails. Known non-Archimedean failures of orthomodularity also warn against
unconditional projection theorems \cite{AN,ANS}. The theorem above gives a
particularly rigid witness: it is the closed range of a bounded, coercive,
self-adjoint injection and retains a continuum-sized algebraic deficit.

\subsection{A continuous functional outside inner-product duality}
\begin{proposition}[Explicit failure of valuation Riesz representation]\label{prop:riesz}
There is a nonzero $\K$-linear, valuation-continuous functional
$F:\HH\to\K$ that vanishes on $M=C\HH$ but is not of the form
$F(x)=\ip{x}{z}$ for any $z\in\HH$.
\end{proposition}
\begin{proof}
Choose an algebraic complement $W$ of $\Ran D^2$ in $H$, and take the
projection $\Pi$ in \cref{thm:defect} for $S=D^2$, $E=s^2I$.
Its coefficients have nonnegative exponents: $U$ and $S$ are constant,
$E$ has positive order, and $(I+UE)^{-1}$ has nonnegative order.
Consequently $v(\Pi x)\ge v(x)$.
Choose a nonzero complex linear functional $\psi:W\to\C$, and extend it
coefficientwise to $\psi_\Gamma:W((t^\Gamma))\to\K$. Set
$F=\psi_\Gamma\Pi$. Coefficient extraction gives
$v(Fx)\ge v(x)$, so $F$ is valuation-continuous. It is nonzero because
$\Pi$ is the identity on $W((t^\Gamma))$, and it kills $M$.
If $F(x)=\ip{x}{z}$, then vanishing on $M$ implies $z\in M^\perp=0$,
which would force $F=0$.
\end{proof}

The algebraic maps $U$ and $\psi$ need not be continuous for the ordinary
Hilbert norm. That does not prevent the displayed Hahn maps from being
valuation-continuous. Confusing those topologies would incorrectly apply
an ordinary closed graph or Riesz argument. Conversely, this proposition
does not contradict \cref{thm:adjoint}: the functional just constructed
has no adjoint represented by a Hahn vector.

\section{When does an operator have a least field-valued norm?}\label{sec:norm}

\subsection{Bounds versus the existence of their least element}
For $A\in\Acal$, let
\[
 \mathcal U(A)=\{M\in\KR_{\ge0}:
             \norm{Ax}\le M\norm{x}\text{ for all }x\in\HH\}.
\]
\Cref{cor:autobounded} makes this set nonempty. By a field-valued
\emph{operator norm} we mean a least element of this set, not an arbitrary
chosen bound and not the valuation norm of the operator.

\begin{lemma}[Ordinary contraction bounds survive]\label{lem:contract}
If $T\in\B$ and $n=\norm{T}_H$, then
$\norm{Tx}\le n\norm{x}$ for every $x\in\HH$.
For nonzero $x$ with leading coefficient $u$, the finite scalar
$\norm{Tx}/\norm{x}$ has standard part
\[
 \st\!\left(\frac{\norm{Tx}}{\norm{x}}\right)
       =\frac{\norm{Tu}_H}{\norm{u}_H}.
\]
The right side is zero if $Tu=0$.
\end{lemma}
\begin{proof}
The ordinary positive operator $n^2I-T^*T$ has a bounded positive square
root $R$ by ordinary Hilbert spectral theory. The coefficientwise identity
$R^*R=n^2I-T^*T$ implies
$n^2\norm{x}^2-\norm{Tx}^2=\norm{Rx}^2\ge0$.
For the standard part, use the leading norm formula in
\cref{prop:inner}. If $Tu=0$, the output has strictly larger valuation
or is zero, so the quotient has residue zero.
\end{proof}

\begin{theorem}[Norm-attainment criterion]\label{thm:norm}
For an ordinary bounded operator $T$ on nonzero $H$, the set
$\mathcal U(T)$ has a least element if and only if $T$ attains its ordinary
operator norm on a unit vector. When it does, the least element is
$\norm{T}_H$ and is attained by a constant unit vector in $\HH$.
\end{theorem}
\begin{proof}
Put $n=\norm{T}_H$. If $T$ attains $n$, the constant maximizing vector
forces every upper bound to be at least $n$, while \cref{lem:contract}
proves that $n$ is a bound. This includes $T=0$, which attains its norm
on every unit vector.

Now suppose $T$ does not attain its norm; then $n>0$. For each nonzero
$x\in\HH$, \cref{lem:contract} gives
\[
 \st\!\left(\frac{\norm{Tx}}{\norm{x}}\right)<n.
\]
Therefore every positive finite $M$ with $\st(M)=n$ is an upper bound,
including $n-\varepsilon$ for every positive infinitesimal $\varepsilon$.
Every finite $M$ with residue greater than $n$ and every positive infinite
$M$ is a bound as well.

No finite nonnegative $M$ with residue less than $n$ can be a bound:
choose an ordinary unit vector $u$ with
$\norm{Tu}_H>\st(M)$ using the definition of the ordinary supremum. Its
constant extension violates that bound. We have thus described all upper
bounds. A bound with residue greater than $n$, or an infinite bound, can
be decreased to $n$. A bound with residue $n$ can be decreased by any
fixed positive infinitesimal and remains positive with the same residue.
There is no least upper bound.
\end{proof}

\begin{example}[A bounded self-adjoint operator without a field-valued norm]
Let $Te_n=(1-1/(n+1))e_n$ on $\ell^2(\N_{\ge1})$. Its ordinary norm is
one but is not attained. Indeed, for every nonzero $u$,
\[
 \norm{u}_H^2-\norm{Tu}_H^2
 =\sum_{n\ge1}\left(1-\left(1-\frac1{n+1}\right)^2\right)|u_n|^2>0,
\]
where this is an ordinary convergent real sum. The diagonal values approach
one, proving the supremum is one. Hence its Hahn extension has many
field-valued bounds, including $1-s$ and $1-2s$, but no least one.
\end{example}

\subsection{The two mechanisms are independent}
The operator $D^2+s^2I$ in \cref{thm:coercive} has an attained least upper
norm bound $1+s^2$ but is not invertible. The operator in the preceding
example fails norm attainment for a different reason: every individual
Hahn vector has a strictly submaximal \emph{ordinary leading ratio}, so
infinitesimally smaller bounds cannot be distinguished by any one vector.
Neither phenomenon should be used as an explanation for the other.

\section{Finite sections, exact checks, and reproducibility}\label{sec:checks}

\subsection{Why finite invertibility is misleading}
Let $D_N=\operatorname{diag}(1,1/2,\ldots,1/N)$ and
$C_N=D_N^2+s^2I$. Every $C_N$ is invertible over $\K$. For the truncated
right-hand side $f_N=(n^{-2})_{1\le n\le N}$,
\begin{equation}\label{eq:finitesolution}
 (C_N^{-1}f_N)_n
 =\frac{1}{1+n^2s^2}
 =\sum_{k\ge0}(-1)^k n^{2k}s^{2k}.
\end{equation}
The constant coefficient of this finite solution is the $N$-vector of
ones, whose ordinary Hilbert norm is $\sqrt N$. No single ordinary
$\ell^2$ vector can serve as the stable constant coefficient as
$N\to\infty$. This directly identifies the obstruction missed by every
finite matrix determinant.

For $A_N=D_N-isI$, all diagonal resolvents likewise exist, with
\[
 (n^{-1}-is)^{-1}
   =\sum_{k\ge0}i^k n^{k+1}s^k.
\]
The leading inverse coefficient has ordinary norm $N$. Infinite extension
would require the unbounded ordinary operator
$\operatorname{diag}(n)$ as a bounded coefficient, contradicting
\cref{thm:adjoint}. More strongly, particular constant data fail the
vector-space condition of \cref{thm:smooth}; thus the problem is actual
nonsurjectivity, not just exclusion from a preferred inverse algebra.

\subsection{The accompanying exact checks}
The program \texttt{code/verify.py} uses exact rational and symbolic
arithmetic. It checks finite-support convolution and adjoints, including a
lexicographically ordered two-dimensional exponent group; truncated
noncommutative Neumann identities; the isolated-point resolvent formula;
the algebraic left-inverse and complement projection identities in a
finite rectangular analogue; the finite-section formulas above; and the
backward-shift residual. The recorded run passed 218 exact assertions in six groups; its output is
included in \texttt{data/verification.txt}.

The finite rectangular analogue is deliberately identified as such. An
injective nonsurjective endomorphism cannot exist on a finite-dimensional
space, so a finite square-matrix test cannot validate the infinite range
defect. Nor can finite tests establish Baire category, arbitrary Hahn
support well-ordering, an infinite cokernel dimension, or the spectral
formula. Those assertions rest on the proofs, not on the program.

The LaTeX source contains its bibliography and needs no external images,
fonts, or bibliography database. A standard build is
\begin{verbatim}
pdflatex -interaction=nonstopmode -halt-on-error article.tex
pdflatex -interaction=nonstopmode -halt-on-error article.tex
pdflatex -interaction=nonstopmode -halt-on-error article.tex
python code/verify.py
\end{verbatim}
The Python script needs SymPy. The source is not Lean code, and no
proof-assistant certification is claimed.

\section{Interpretation and remaining boundaries}\label{sec:boundaries}

\subsection{What has been settled in this construction}
The questions raised by this particular scalar-extension model now have
sharp answers. Adjointability is not an opaque class condition: it is
exactly global Hahn support with bounded ordinary coefficients. For
constant normal operators, infinitesimal resolvent behavior is completely
classified by the ordinary spectral accumulation set. Infinitesimal
perturbations of injective leading operators cannot remove any ordinary
range defect. The leading ordinary coefficient gives a direct and robust
witness, even where every finite section is invertible.

These statements distinguish spectral membership from eigenvalue
existence, compactness in the original Hilbert topology from compactness in
the valuation topology, upper norm bounds from a least such bound, and
metric completeness from orthogonal completeness. None of those pairs can
be silently identified after extending the scalars.

\subsection{What would require a different theorem}
The formula \eqref{eq:exactspectrum} concerns \emph{constant normal}
operators. General Hahn operators with multiple coefficient scales can
have more complicated spectra; \cref{ex:backshift} already shows why the
point-spectrum rigidity fails without normality. The perturbation theorem
handles injective leading coefficients and proves a large class of
obstructions, but is not a classification of all spectra of all Hahn
operator series.

One could change the vector construction to allow coefficientwise objects
outside $H$, choose a different sequence-space completion, or impose a
stronger orthogonal-completeness axiom. Such changes must be specified and
analyzed rather than treated as harmless notational variations. Enlarging
only the exponent group is not enough, by \cref{cor:extension}.

The full surreal class offers additional scales but is not a set-sized
coefficient field in ZFC. Our proofs avoid class-indexed summations and
class-sized choices by fixing the workspace at the outset. The surcomplex
interpretation \eqref{eq:surreal-embedding} supplies genuine infinitesimal
and infinite scalars while keeping every construction set-sized.

\subsection{Why this matters beyond a counterexample}
In finite dimension, replacing $\R/\C$ by a real-closed field and its
complexification preserves much of spectral algebra. In infinite dimension,
that replacement cannot by itself preserve the analytic role of Hilbert
space. An infinitesimal imaginary displacement may fail to produce a
resolvent, even when it gives a strict lower norm bound. An infinitesimal
positive regularizer may leave a continuum of unsolved linear equations.
These are precise obstructions to using the model as an automatic
replacement for complex Hilbert space, not a rejection of surcomplex
analysis in every possible formulation.

The constructive side is equally useful. The automatic adjoint theorem
gives a well-defined algebra for exact symbolic manipulation. The inverse
formulas identify every ordinary normal resolvent that actually exists.
The inverse-smooth criterion identifies exactly which constant data remain
solvable at a monomial infinitesimal shift. The defect splitting gives
explicit noncanonical projections and an exact algebraic quotient. Together
these form a workable theory with both positive structure and sharp limits.

\appendix
\section{Repository and literature audit}\label{app:audit}

\subsection{Repository snapshot}
The coverage comparison is pinned to
\begin{center}\small\texttt{\commit}.\end{center}
We inspected the repository tree, the project \texttt{README.md}, the full
\texttt{docs/README.md} catalogue, the opening 250 lines of
\texttt{docs/surcomplex/spectral-theory/article.tex}, and the opening 230
lines of the differential-equations article. This was a targeted scope
audit, not a claim to have reread every archived source manuscript or
verified every existing Lean file. The spectral opening explicitly states
its finite-dimensional scope, while the catalogue maps the already-covered
analytic, differential, Berkovich, foundational and computational strands.

\subsection{Primary-source comparison and date boundary}
The literature search was carried out through September 22, 2026. The
publisher records and accessible abstracts for \cite{ANS,Ishizuka}, the
author-provided text associated with \cite{AN}, and the arXiv records for
\cite{Boasso,Miabey} were used to locate nearby topics. Ishizuka's paper
appears in volume 39(2), labeled 2024, with the publisher recording online
publication on January 3, 2025. Both dates are recorded in the bibliography
to avoid silently conflating the issue year with the online date.

The comparison supports the modest statement that the exact theorem
package here was not identified in the inspected sources. It does not
establish that no earlier work contains any equivalent result. In
particular, formal support inversion, failure of non-Archimedean
orthogonality, and ordinary resolvent expansions have substantial
pre-existing literatures.

\section{Dependency and proof-obligation ledger}\label{app:dependencies}
\begin{table}[ht]
\centering\small
\begin{tabularx}{\textwidth}{@{}p{.27\textwidth}YY@{}}
\toprule
Result & Inputs & Independent restriction\\
\midrule
Positive inner product & Finite Hahn convolution; positive square roots & Divisible value group\\
Automatic adjoint structure & Ordinary closed graph; Baire; Hahn support & Everywhere-defined adjoints\\
Defect splitting & Injective leading term; Neumann support lemma & No topology or commutation needed\\
Normal spectral formula & Defect splitting; ordinary normal spectral theorem & Constant normal operator\\
New eigenvalue exclusion & Orthogonal eigenspace projection & Normality is essential\\
Continuum cokernel & Leading-coefficient obstruction; power-series divergence test & Ordinary $\ell^2$ coefficients\\
Closed coercive range & Defect valuation equality; metric completeness & Rank-one value group\\
Least norm-bound criterion & Ordinary positive square root; leading norm ratio & Constant bounded operator\\
\bottomrule
\end{tabularx}
\caption{Where each hypothesis enters; ``complete'' does not replace ``orthogonally complete.''}
\end{table}

All general statements in the main theorem package have proofs in the text,
relative to the named classical inputs. No computational observation is
used as a premise of those proofs. The most consequential steps to audit
independently are the Baire uniform-support argument, the complement
projection in \cref{thm:defect}, and the accumulation-point case in
\cref{thm:spectrum}; their details are given without an appeal to a
numerical experiment.

\clearpage
\phantomsection
\addcontentsline{toc}{section}{References}
\begin{thebibliography}{99}
\raggedright
\bibitem{RepoMap}
V. Reshetnikov, \emph{Surreal repository: documentation catalogue},
\texttt{docs/README.md}, revision \texttt{\commit}, inspected September 22, 2026.
\url{https://github.com/VladimirReshetnikov/Surreal/blob/608dd23c0539fce73723843949ee627641c27b30/docs/README.md}.

\bibitem{RepoSpectral}
\emph{Surcomplex Spectral Theory: Singular Values, Infinitesimal Rank, and
Multiscale Stability}, AI-assisted research exposition in the Surreal
repository, \texttt{docs/surcomplex/spectral-theory/article.tex}, revision
\texttt{\commit}, September 2026.
\url{https://github.com/VladimirReshetnikov/Surreal/blob/608dd23c0539fce73723843949ee627641c27b30/docs/surcomplex/spectral-theory/article.tex}.

\bibitem{Gonshor}
H. Gonshor, \emph{An Introduction to the Theory of Surreal Numbers},
London Mathematical Society Lecture Note Series 110, Cambridge University
Press, 1986.
\url{https://www.cambridge.org/core/books/an-introduction-to-the-theory-of-surreal-numbers/312AE504A3E88E804054BFB390446374}.

\bibitem{Neumann}
B. H. Neumann, On ordered division rings,
\emph{Transactions of the American Mathematical Society} \textbf{66}
(1949), 202--252.

\bibitem{Douglas}
R. G. Douglas, \emph{Banach Algebra Techniques in Operator Theory},
second edition, Graduate Texts in Mathematics 179, Springer, 1998.
\url{https://doi.org/10.1007/978-1-4612-1656-8}.

\bibitem{AN}
J. Aguayo and M. Nova, Non-Archimedean Hilbert like spaces,
\emph{Bulletin of the Belgian Mathematical Society -- Simon Stevin}
\textbf{14} (2007), no.~5, 787--797.
\url{https://doi.org/10.36045/bbms/1197908895}.

\bibitem{ANS}
J. Aguayo, M. Nova and K. Shamseddine,
Characterization of compact and self-adjoint operators on free Banach
spaces of countable type over the complex Levi-Civita field,
\emph{Journal of Mathematical Physics} \textbf{54} (2013), 023503.
\url{https://doi.org/10.1063/1.4789541}.

\bibitem{Ishizuka}
K. Ishizuka, A spectral theorem for a non-Archimedean valued field whose
residue field is formally real,
\emph{Extracta Mathematicae} \textbf{39} (2024), no.~2, 235--254;
publisher online date January 3, 2025.
\url{https://doi.org/10.17398/2605-5686.39.2.235}.

\bibitem{Boasso}
E. Boasso, The Drazin spectrum in Banach algebras,
in \emph{An Operator Theory Summer: Timisoara, June 29--July 4, 2010},
International Book Series of Mathematical Texts, Theta Foundation,
Bucharest, 2012, 21--28; arXiv:1309.5025, 2013.
\url{https://arxiv.org/abs/1309.5025}.

\bibitem{Miabey}
T. Miabey, Spectral and Essential Spectral Analysis of Finite-Rank
Perturbations of Unbounded Diagonal Operators on Non-Archimedean Hilbert
Spaces, arXiv:2606.06783v1, June 5, 2026, preprint.
\url{https://arxiv.org/abs/2606.06783}.
\end{thebibliography}
\end{document}
